\documentclass[aspectratio=169,10pt]{beamer}
% Self-contained replacement for the shared CS 4880 Beamer preamble.
% Uses the same macros, naming pattern, concept-lead / keyidea style as Lecture 31.
\usetheme{default}
\usefonttheme{professionalfonts}
\definecolor{CalPolyGreen}{HTML}{154734}
\definecolor{DeckSoft}{HTML}{EAF2ED}
\definecolor{DeckInk}{HTML}{263A32}
\definecolor{DeckMuted}{HTML}{63756B}
\definecolor{DeckGold}{HTML}{B88A2B}
\definecolor{DeckMint}{HTML}{168F85}
\setbeamercolor{normal text}{fg=DeckInk,bg=white}
\setbeamercolor{frametitle}{fg=CalPolyGreen,bg=white}
\setbeamercolor{structure}{fg=CalPolyGreen}
\setbeamercolor{block title}{fg=white,bg=CalPolyGreen}
\setbeamercolor{block body}{fg=DeckInk,bg=DeckSoft}
\setbeamercolor{itemize item}{fg=CalPolyGreen}
\setbeamercolor{itemize subitem}{fg=CalPolyGreen}
\setbeamercolor{enumerate item}{fg=CalPolyGreen}
\setbeamertemplate{navigation symbols}{}
\setbeamertemplate{itemize item}[circle]
\setbeamertemplate{itemize subitem}[triangle]
\setbeamerfont{frametitle}{size=\large,series=\bfseries}
\setbeamertemplate{frametitle}{%
  \nointerlineskip\vspace*{.08cm}\hspace{.34cm}\insertframetitle\par
  \hspace{.34cm}{\color{CalPolyGreen}\rule{.85\paperwidth}{.7pt}}\vspace{.04cm}
}
\setbeamertemplate{footline}{%
  \leavevmode\hbox{\begin{beamercolorbox}[wd=.83\paperwidth,ht=2.8ex,dp=1.2ex,leftskip=.43in]{normal text}%
  {\scriptsize\color{DeckMuted}\insertshortauthor\quad\insertshorttitle}%
  \end{beamercolorbox}%
  \begin{beamercolorbox}[wd=.17\paperwidth,ht=2.8ex,dp=1.2ex,rightskip=.43in plus1fil]{normal text}%
  {\scriptsize\color{DeckMuted}\hfill\insertframenumber/\inserttotalframenumber}%
  \end{beamercolorbox}}%
}
\setbeamertemplate{headline}{}
\RequirePackage{tikz}
\newcommand{\coursenumber}[1]{\def\csCourseNumber{#1}}
\newcommand{\coursetitle}[1]{\def\csCourseTitle{#1}}
\newcommand{\lecturenumber}[1]{\def\csLectureNumber{#1}}
\newcommand{\lecturetitle}[1]{\def\csLectureTitle{#1}\title[Natural Language Processing]{#1}}
\newcommand{\lecturedate}[1]{\def\csLectureDate{#1}}
\newcommand{\lecturedisplaytitle}[1]{\def\csDisplayTitle{#1}}
\newcommand{\courseunit}[1]{\def\csCourseUnit{#1}}
\newcommand{\makedecktitle}{%
\begin{frame}[plain]
\begin{tikzpicture}[remember picture,overlay]
\fill[CalPolyGreen] (current page.north west) rectangle ([xshift=0.48in]current page.south west);
\end{tikzpicture}
\vspace*{.48in}\hspace*{.60in}%
\begin{minipage}{.75\paperwidth}
{\large\bfseries\color{DeckMuted}\csCourseNumber\quad\csCourseTitle\par}\vspace{.30in}
{\fontsize{24}{29}\selectfont\bfseries\color{CalPolyGreen}\raggedright\csDisplayTitle\par}\vspace{.20in}
{\Large\color{DeckInk}Lecture \csLectureNumber\par}\vspace{.28in}
{\normalsize\color{DeckMuted}\csCourseUnit\quad\textbullet\quad\csLectureDate\par}
\end{minipage}
\end{frame}}
\newenvironment{keyidea}{%
\par\vspace{3pt}\begin{beamercolorbox}[rounded=true,shadow=false,sep=5pt,wd=\linewidth]{block body}%
\color{CalPolyGreen}\textbf{Key idea.}\,\color{DeckInk}%
}{\end{beamercolorbox}}
\author[CS 4880]{CS 4880 Artificial Intelligence}

\usepackage[T1]{fontenc}
\usepackage{lmodern}
\usepackage{amsmath,amssymb,booktabs,array}
\usepackage{graphicx}
\graphicspath{{../assets/lecture-35/}}
\pdfstringdefDisableCommands{\def\translate#1{#1}}
% CS4880 Lecture 35, November 20, 2026. About 50 minutes.
% Arc: problem & history 5; inputs 8; single/multi-head attention 19;
% transformer blocks/families 8; training/generation 7; limits & exit 3.
% This lesson combines the provided 24_480 and 25_480 lecture PDFs.
% All diagrams are generated by generate-figures.py and labelled schematic
% unless their numerical values are explicitly derived and reproducible.
% Keep the mathematical distinction between 2017 encoder-decoder Transformers
% and causal decoder-only GPT models explicit.
\newcommand{\lecturevisual}[2][2.15in]{%
  \par{\centering\includegraphics[width=.98\linewidth,height=#1,keepaspectratio]{#2}\par}%
}
\newcommand{\conceptlead}[1]{{\small\bfseries\color{CalPolyGreen}#1\par}}
\newenvironment{calculation}[1][Mathematical description]{%
  \par\begingroup\setbeamercolor{calculation body}{bg=DeckSoft,fg=DeckInk}%
  \begin{beamercolorbox}[sep=4pt,wd=\linewidth]{calculation body}%
  \footnotesize{\bfseries\color{DeckMuted}#1\par}%
  \setlength{\jot}{0pt}\setlength{\abovedisplayskip}{3pt}\setlength{\belowdisplayskip}{3pt}%
}{\end{beamercolorbox}\endgroup}
\newcommand{\microhead}[1]{{\footnotesize\bfseries\color{CalPolyGreen}#1\par}}
\setbeamersize{text margin left=.42in,text margin right=.42in}
\coursenumber{CS 4880}
\coursetitle{Artificial Intelligence}
\lecturenumber{35}
\lecturetitle{Transformers: From Tokens to Attention}
\lecturedisplaytitle{Transformers:\\[3pt]From Tokens to Attention}
\lecturedate{November 20, 2026}
\courseunit{Natural Language Processing}
\begin{document}
\makedecktitle

\section{The Problem That Attention Addresses}

\begin{frame}{Last Lecture Left Us With One Bottleneck}
\small
\conceptlead{A seq2seq encoder had to summarize an entire source sentence into a fixed-size representation.}
\begin{columns}[T,onlytextwidth]
\begin{column}{.60\textwidth}\lecturevisual[2.22in]{seq2seq_bottleneck.png}\end{column}
\begin{column}{.37\textwidth}
\begin{itemize}
\item The decoder may need one particular source word much later.
\item A recurrent hidden state must carry that information forward.
\item Longer sequences can make learning these dependencies difficult.
\end{itemize}
\end{column}
\end{columns}
\begin{keyidea}What if a token could retrieve relevant information directly, instead of asking one vector to remember everything?\end{keyidea}
\end{frame}

\begin{frame}{A Small Historical Pivot: 2014 to 2017}
\small
\conceptlead{Attention was developed for sequence models before the Transformer.}
\lecturevisual[1.54in]{history_to_transformer.png}
\begin{itemize}
\item \textbf{2014:} attention-enhanced neural machine translation (Bahdanau, Cho, Bengio) lets a decoder consult source states.
\item \textbf{2017:} Vaswani et al., \emph{Attention Is All You Need}, remove recurrence from the core sequence model.
\item \textbf{Why it mattered:} better parallelization of token computations \emph{during training}, plus direct pathways between positions.
\end{itemize}
\end{frame}

\begin{frame}{Today's Question: How Does a Token Use Context?}
\small
\begin{block}{By the end of the lecture}
\begin{enumerate}
\item Trace text through tokens, embeddings, positions, transformer layers, and logits.
\item Compute one row of scaled dot-product attention with numbers.
\item Explain multiple heads, a causal mask, residual connections, and feed-forward layers.
\item Distinguish encoder-only, decoder-only, and encoder--decoder models.
\item Explain how next-token training differs from generation at inference time.
\end{enumerate}
\end{block}
\begin{keyidea}The central operation is a \emph{learned, context-dependent weighted combination of value vectors}.\end{keyidea}
\end{frame}

\section{Turning Text Into Inputs}

\begin{frame}{First See the Entire Language-Model Pipeline}
\small
\lecturevisual[2.55in]{llm_pipeline.png}
\begin{itemize}
\item \textbf{Input processing:} a tokenizer converts text into discrete token IDs.
\item \textbf{Representation:} an embedding table and position information produce vectors.
\item \textbf{Model:} stacked transformer blocks compute context-aware features.
\item \textbf{Output:} a head produces logits over tokens; a decoding strategy chooses the continuation.
\end{itemize}
\end{frame}

\begin{frame}{Step 1: Tokenization Is a Design Choice}
\small
\conceptlead{Text is not automatically split into words. Subword vocabularies reuse frequent pieces.}
\begin{columns}[T,onlytextwidth]
\begin{column}{.60\textwidth}\lecturevisual[2.33in]{subword_example.png}\end{column}
\begin{column}{.37\textwidth}
\begin{itemize}
\item Vocabulary entries can be words, parts of words, symbols, or bytes.
\item Actual splits vary between tokenizers: BPE and SentencePiece are examples.
\item Different tokenizations change sequence length and available context.
\end{itemize}
\end{column}
\end{columns}
\begin{keyidea}A token is a unit in a particular model's vocabulary---not necessarily a full word.\end{keyidea}
\end{frame}

\begin{frame}{Step 2: A Token ID Selects an Embedding Vector}
\small
\conceptlead{The ID itself is just an index; the associated vector is learned.}
\begin{columns}[T,onlytextwidth]
\begin{column}{.58\textwidth}\lecturevisual[2.45in]{embedding_lookup.png}\end{column}
\begin{column}{.39\textwidth}
\begin{calculation}[Embedding lookup]
\[E\in\mathbb R^{|\mathcal V|\times d},\qquad x_i=E[t_i].\]
$|\mathcal V|$ is vocabulary size, $d$ is model width, and $t_i$ is an integer token ID.
\end{calculation}
\begin{itemize}
\item Learned embeddings begin as parameters, not dictionary definitions.
\item Similarity can emerge through training, but depends on the objective.
\end{itemize}
\end{column}
\end{columns}
\end{frame}

\begin{frame}{Step 3: Without Position, Word Order Can Disappear}
\small
\lecturevisual[2.23in]{word_order.png}
\begin{itemize}
\item Bare self-attention, without any positional signal, is equivariant to reordering its input tokens.
\item That is a problem for syntax: \emph{``dog bites human''} differs from \emph{``human bites dog''}.
\item Add positional information so the model can represent ordering and distances.
\end{itemize}
\end{frame}

\begin{frame}{A Position Vector Changes Each Token's Input}
\small
\conceptlead{The original 2017 Transformer used sinusoidal positions; some later models learn position embeddings.}
\begin{columns}[T,onlytextwidth]
\begin{column}{.56\textwidth}\lecturevisual[2.25in]{sinusoidal_position.png}\end{column}
\begin{column}{.41\textwidth}
\begin{calculation}[Simplified additive approach]
\[h_i^{(0)}=E[t_i]+p_i.\]
The two vectors share the same width $d$.
\end{calculation}
\begin{itemize}
\item Original sinusoidal positions use different sine/cosine frequencies.
\item Other models incorporate position through relative or rotary methods.
\end{itemize}
\end{column}
\end{columns}
\end{frame}

\begin{frame}{Static Embedding Versus Contextual Representation}
\small
\conceptlead{The token ``bank'' may start with the same lexical embedding, but its meaning depends on its neighbors.}
\lecturevisual[1.40in]{contextual_bank.png}
\begin{itemize}
\item Word2vec generally learns a fixed vector for a word type.
\item Transformer layers repeatedly update each position based on surrounding tokens.
\end{itemize}
\begin{keyidea}A word's input embedding is a starting point; a transformer produces representations \emph{conditioned on this sentence}.\end{keyidea}
\end{frame}

\section{One Attention Head, Step by Step}

\begin{frame}{Attention: Which Other Tokens Matter Here?}
\small
\conceptlead{A token can blend information from other positions, with different weights for different contexts.}
\lecturevisual[2.65in]{attention_intuition.png}
\begin{itemize}
\item The diagram suggests a plausible link between \emph{it} and \emph{glass}; actual attention weights are learned, not supplied by grammar.
\item Attention is a general vector operation, not an explicit symbolic pronoun resolver.
\end{itemize}
\end{frame}

\begin{frame}{From Each Input Vector, Make Three New Vectors}
\small
\lecturevisual[2.05in]{qkv_projections.png}
\begin{calculation}[One layer, one head]
\[q_i=x_i W_Q,\qquad k_i=x_i W_K,\qquad v_i=x_i W_V.\]
$W_Q,W_K,W_V$ are separate learned matrices (bias terms omitted here).
\end{calculation}
\end{frame}

\begin{frame}{Queries, Keys, and Values Have Distinct Jobs}
\small
\begin{columns}[T,onlytextwidth]
\begin{column}{.31\textwidth}
\begin{block}{Query $q_i$}
What information does position $i$ \emph{seek} from the sequence?
\end{block}
\end{column}
\begin{column}{.31\textwidth}
\begin{block}{Key $k_j$}
How well does position $j$ \emph{match} that query?
\end{block}
\end{column}
\begin{column}{.35\textwidth}
\begin{block}{Value $v_j$}
What information does position $j$ \emph{contribute} if attended to?
\end{block}
\end{column}
\end{columns}
\begin{keyidea}Queries and keys determine the weights; values are what gets combined. The three are learned projections, not human-written tags.\end{keyidea}
\end{frame}

\begin{frame}{Step 1: Compare a Query Against Every Key}
\small
\conceptlead{The dot product supplies one compatibility score for each candidate position.}
\begin{calculation}[Scaled dot-product compatibility]
\[s_{ij}=\frac{q_i\cdot k_j}{\sqrt{d_k}}.\]
$q_i,k_j\in\mathbb R^{d_k}$. A larger score contributes more weight \emph{after softmax}, all else equal.
\end{calculation}
\begin{itemize}
\item The dot product depends on vector directions \emph{and magnitudes}; it is not automatically cosine similarity.
\item Division by $\sqrt{d_k}$ keeps typical score magnitudes manageable as key width grows.
\item Scores are \emph{not} yet probabilities; they can be negative or greater than one.
\end{itemize}
\end{frame}

\begin{frame}{Step 2: Turn the Scores Into Weights}
\small
\begin{calculation}[Row-wise softmax]
\[\alpha_{ij}=\frac{\exp(s_{ij})}{\sum_{\ell\in\mathcal A_i}\exp(s_{i\ell})},\qquad \sum_{j\in\mathcal A_i}\alpha_{ij}=1.\]
$\mathcal A_i$ is the set of key positions allowed by the attention mask.
\end{calculation}
\begin{itemize}
\item Each query has its \emph{own} distribution over allowed keys.
\item Exponentials make weights positive; normalization makes them sum to one.
\item A mask may remove positions before softmax.
\end{itemize}
\begin{keyidea}Softmax converts unnormalized matching scores into a convex set of weights over value vectors.\end{keyidea}
\end{frame}

\begin{frame}{Worked Example: Three Keys, One Query}
\small
\conceptlead{We will compute this entire attention row by hand with $d_k=2$.}
\lecturevisual[1.90in]{scores_to_softmax.png}
\begin{calculation}[Exact toy setup]
\[q=(1,0),\quad k_1=(1,0),\ k_2=(0,1),\ k_3=(1,1).\]
\[s=(1/\sqrt2,\ 0,\ 1/\sqrt2)\approx(0.707,0,0.707).\]
\end{calculation}
\end{frame}

\begin{frame}{Worked Example: Check the Softmax}
\small
\begin{calculation}[Calculate normalized weights]
\[e^{0.7071}\approx2.028,\quad e^0=1,\quad Z=2.028+1+2.028\approx5.056.\]
\[\alpha=\left(\frac{2.028}{5.056},\frac{1}{5.056},\frac{2.028}{5.056}\right)\approx(0.4011,0.1978,0.4011).\]
\end{calculation}
\begin{itemize}
\item The first and third keys match equally well; both receive about $40\%$.
\item The second key still contributes nearly $20\%$.
\item Changing the query changes the row of weights.
\end{itemize}
\end{frame}

\begin{frame}{Step 3: Use Those Weights to Combine Values}
\small
\lecturevisual[1.67in]{value_weighted_sum.png}
\begin{calculation}[Output vector for the query position]
\[v_1=(2,0),\quad v_2=(0,2),\quad v_3=(1,1).\]
\[z=0.4011v_1+0.1978v_2+0.4011v_3\approx(1.2033,0.7967).\]
\end{calculation}
\begin{keyidea}This $z$ is a new feature vector, not a selected word, sentence, or class label.\end{keyidea}
\end{frame}

\begin{frame}{All Positions At Once: Matrix Attention}
\small
\conceptlead{Compute every query's weights in parallel using matrix multiplication.}
\begin{calculation}[Scaled dot-product attention]
\[\boxed{\operatorname{Attention}(Q,K,V)=\operatorname{softmax}\!\left(\frac{QK^\top}{\sqrt{d_k}}+M\right)V}\]
\[Q,K\in\mathbb R^{n\times d_k},\quad V\in\mathbb R^{n\times d_v},\quad \text{output}\in\mathbb R^{n\times d_v}.\]
$M$ is zero for permitted pairs and $-\infty$ for forbidden pairs (conceptually); softmax acts \emph{across each row}.
\end{calculation}
\begin{itemize}
\item $QK^\top$ is $n\times n$: each row is a query position, each column a key position.
\item The result multiplies $V$ to produce one contextualized vector per query position.
\end{itemize}
\end{frame}

\begin{frame}{Read an Attention Matrix Carefully}
\small
\begin{columns}[T,onlytextwidth]
\begin{column}{.57\textwidth}\lecturevisual[2.62in]{toy_attention_heatmap.png}\end{column}
\begin{column}{.40\textwidth}
\begin{itemize}
\item A \textbf{row} sums to 1 and answers: where does this query put weight?
\item A \textbf{column} identifies a candidate source position.
\item Different queries have different rows.
\item This heatmap is computed from toy matrices, \emph{not a trained language model}.
\end{itemize}
\end{column}
\end{columns}
\end{frame}

\begin{frame}{A Three-Line NumPy Check}
\footnotesize
\conceptlead{The same calculation can be verified without a neural-network library.}
\begin{block}{Run after \texttt{import numpy as np}}
\ttfamily
q = np.array([1., 0.])\par
K = np.array([[1.,0.], [0.,1.], [1.,1.]])\par
V = np.array([[2.,0.], [0.,2.], [1.,1.]])\par
s = K @ q / np.sqrt(2)\par
w = np.exp(s - s.max()); w /= w.sum()\par
print(np.round(w, 4))  \# [0.4011 0.1978 0.4011]\par
print(np.round(w @ V, 4)) \# [1.2033 0.7967]
\end{block}
\begin{keyidea}Attention is concrete linear algebra: projections, dot products, softmax, and weighted sums.\end{keyidea}
\end{frame}

\begin{frame}{Why Is It Called \emph{Self}-Attention?}
\small
\begin{columns}[T,onlytextwidth]
\begin{column}{.48\textwidth}
\begin{block}{Self-attention}
\begin{itemize}
\item $Q$, $K$, and $V$ are projected from the \emph{same} sequence.
\item Positions communicate within the sequence (subject to masking).
\end{itemize}
\end{block}
\end{column}
\begin{column}{.48\textwidth}
\begin{block}{Cross-attention}
\begin{itemize}
\item Queries come from one sequence; keys/values come from another.
\item Example: translation decoder consulting encoder outputs.
\end{itemize}
\end{block}
\end{column}
\end{columns}
\begin{keyidea}``Attention'' names a computation; \emph{where Q, K, and V originate} determines the attention type.\end{keyidea}
\end{frame}

\begin{frame}{One Head Is Limited to One Set of Projections}
\small
\conceptlead{The model need not commit to one query/key/value projection.}
\begin{itemize}
\item Different linguistic dependencies may require different feature comparisons.
\item With \textbf{multi-head attention}, the model computes several separate attention operations.
\item The head outputs are concatenated, then passed through a learned output projection.
\end{itemize}
\begin{calculation}[Multiple heads]
\[H_j=\operatorname{Attention}(Q_j,K_j,V_j),\qquad
\operatorname{MHA}(X)=\operatorname{Concat}(H_1,\ldots,H_h)W_O.\]
\end{calculation}
\end{frame}

\begin{frame}{Multi-Head Attention: Several Views of the Same Input}
\small
\lecturevisual[2.17in]{multihead_attention.png}
\begin{itemize}
\item Each head has its \emph{own learned} $W_Q^{(j)},W_K^{(j)},W_V^{(j)}$.
\item Heads can specialize, but no particular syntactic role is guaranteed.
\item The output layer $W_O$ recombines the heads into the model's feature width.
\end{itemize}
\end{frame}

\begin{frame}{Question: Can a Model Peek at the Future?}
\small
\conceptlead{For autoregressive text generation, predicting position $i$ must not use tokens to its right.}
\lecturevisual[1.85in]{causal_mask.png}
\begin{itemize}
\item An \textbf{encoder} commonly uses bidirectional attention: both earlier and later input positions are visible.
\item A \textbf{causal decoder} permits only the prefix up to the current position.
\item Masking enforces this during training, even though all training positions are computed together.
\end{itemize}
\end{frame}

\begin{frame}{Masking Happens Before Softmax}
\small
\begin{calculation}[Causal mask for a four-token sequence]
\[M=\begin{bmatrix}
0&-\infty&-\infty&-\infty\\
0&0&-\infty&-\infty\\
0&0&0&-\infty\\
0&0&0&0
\end{bmatrix}\]
\[A=\operatorname{softmax}\!\left(\frac{QK^\top}{\sqrt{d_k}}+M\right).\]
\end{calculation}
\begin{itemize}
\item For row 1, only key 1 remains: its normalized weight becomes 1.
\item For row 3, keys 1, 2, 3 are allowed; key 4 receives weight 0.
\end{itemize}
\begin{keyidea}Causal masking prevents future-token leakage, but it does \emph{not} eliminate the need to generate tokens sequentially at inference.\end{keyidea}
\end{frame}

\section{Transformer Architectures and Blocks}

\begin{frame}{The 2017 Transformer Was an Encoder--Decoder}
\small
\conceptlead{The original paper was about sequence-to-sequence machine translation.}
\lecturevisual[1.80in]{encoder_decoder.png}
\begin{itemize}
\item Encoder layers build representations of the source sentence.
\item Decoder layers perform causal self-attention over the target prefix and cross-attend to encoder outputs.
\item A modern GPT-style language model is usually \emph{decoder-only}: not the entire original architecture.
\end{itemize}
\end{frame}

\begin{frame}{One Architecture Family, Three Common Designs}
\small
\lecturevisual[1.55in]{model_families.png}
\begin{itemize}
\item \textbf{Encoder-only:} contextual representations with broad input visibility.
\item \textbf{Decoder-only:} causal attention for next-token prediction and generation.
\item \textbf{Encoder--decoder:} encode source, then generate a target sequence with cross-attention.
\end{itemize}
\begin{keyidea}Do not equate ``Transformer'' with ``GPT''; the masking and network layout depend on the task.\end{keyidea}
\end{frame}

\begin{frame}{What Lives Inside One Transformer Block?}
\small
\begin{columns}[T,onlytextwidth]
\begin{column}{.50\textwidth}\lecturevisual[2.78in]{transformer_block.png}\end{column}
\begin{column}{.47\textwidth}
\begin{enumerate}
\item Self-attention mixes information across positions.
\item A residual connection adds the block input; normalization stabilizes representations.
\item A feed-forward network transforms the features of each position.
\item Another residual connection and normalization follow.
\end{enumerate}
\end{column}
\end{columns}
\end{frame}

\begin{frame}{Why Keep Adding the Input Back?}
\small
\conceptlead{Residual connections give each sublayer a direct path for information and gradients.}
\begin{calculation}[Post-norm, as in the original paper]
\[u=\operatorname{LayerNorm}(X+\operatorname{Attention}(X)),\]
\[Y=\operatorname{LayerNorm}(u+\operatorname{FFN}(u)).\]
\end{calculation}
\begin{itemize}
\item The sublayer learns an \emph{update} rather than replacing the entire input signal.
\item Layer normalization normalizes features per token and learns scale/shift parameters.
\item Many later designs use \emph{pre-norm}: normalization is applied before each sublayer.
\end{itemize}
\end{frame}

\begin{frame}{The Feed-Forward Layer Is Shared Across Positions}
\small
\lecturevisual[1.50in]{positionwise_ffn.png}
\begin{calculation}[Original position-wise feed-forward network]
\[\operatorname{FFN}(x)=\max(0,xW_1+b_1)W_2+b_2.\]
\end{calculation}
\begin{itemize}
\item Attention \emph{mixes information between tokens}; the FFN \emph{transforms each token's features}.
\item These weights are shared across positions within a given layer.
\item Later models may replace ReLU with GELU or gated alternatives.
\end{itemize}
\end{frame}

\begin{frame}{Why Stack Several Blocks?}
\small
\conceptlead{One attention operation is a shallow transformation; repeated layers build richer features.}
\begin{enumerate}
\item Start with token embeddings plus position information.
\item A first block forms context-aware vectors from nearby and distant inputs.
\item Subsequent blocks operate on those newly contextualized vectors.
\item The last layer's hidden states feed task-specific heads or token logits.
\end{enumerate}
\begin{keyidea}The stack learns representations end-to-end; it is not a fixed dictionary of grammar rules or human-named concepts.\end{keyidea}
\end{frame}

\section{Predicting, Training, and Generating}

\begin{frame}{From a Context Vector to a Token Distribution}
\small
\conceptlead{The final hidden state is mapped to one logit per vocabulary item.}
\begin{columns}[T,onlytextwidth]
\begin{column}{.59\textwidth}\lecturevisual[2.58in]{softmax_distribution.png}\end{column}
\begin{column}{.38\textwidth}
\begin{calculation}[Illustrative logits]
\[\ell=[0.2,0.3,0.9],\quad p=\operatorname{softmax}(\ell).\]
\[p\approx[.243,.268,.489].\]
\end{calculation}
\begin{itemize}
\item Logits are scores, not probabilities.
\item Softmax assigns mass across the full output vocabulary.
\end{itemize}
\end{column}
\end{columns}
\end{frame}

\begin{frame}{Training: Predict the Next Token at Every Position}
\small
\lecturevisual[2.58in]{teacher_forcing.png}
\begin{itemize}
\item The training sequence provides the prefix tokens; targets are shifted one position forward.
\item A causal mask stops the model from seeing the target token at a future position.
\item All per-position losses can be computed in parallel for a given training sequence.
\end{itemize}
\end{frame}

\begin{frame}{Training Uses Cross-Entropy, Not Human Corrections to Every Sentence}
\small
\begin{calculation}[Token-level negative log-likelihood]
\[\mathcal L=-\frac1T\sum_{t=1}^{T}\log p_\theta(x_t\mid x_{<t}).\]
With a toy true-token probability $p=0.7$, that token's contribution is $-\ln(0.7)\approx0.357$ nats.
\end{calculation}
\begin{itemize}
\item Gradients update embeddings, Q/K/V projections, output layers, and other parameters.
\item Optimization often uses variants of Adam; training data and objective vary by model.
\item Minimizing token loss does not guarantee factual truth or grounded reasoning.
\end{itemize}
\end{frame}

\begin{frame}{Inference: One Token Becomes Part of the Next Input}
\small
\lecturevisual[2.50in]{autoregressive_loop.png}
\begin{itemize}
\item \textbf{Training:} known target sequences allow many token positions to be processed in parallel.
\item \textbf{Generating:} the next token is unknown until the preceding token has been selected.
\item Reusing cached key/value projections can reduce repeated computation, but decoding is still autoregressive.
\end{itemize}
\end{frame}

\begin{frame}{Why Not Always Pick the Largest Probability?}
\small
\begin{columns}[T,onlytextwidth]
\begin{column}{.47\textwidth}
\begin{block}{Greedy decoding}
Choose $\arg\max p_t$ each time.\\[5pt]
Deterministic for a fixed model and input; may produce repetitive or bland outputs.
\end{block}
\end{column}
\begin{column}{.49\textwidth}
\begin{block}{Sampling}
Sample from the distribution, optionally after temperature scaling or top-$k$/top-$p$ truncation.\\[5pt]
Trades diversity and variability against control.
\end{block}
\end{column}
\end{columns}
\begin{calculation}[Temperature before softmax]
\[p_i(T)=\frac{e^{\ell_i/T}}{\sum_j e^{\ell_j/T}},\qquad T>0.\]
$T<1$ sharpens the distribution; $T>1$ flattens it.
\end{calculation}
\end{frame}

\begin{frame}[fragile]{A Compact PyTorch Attention Implementation}
\scriptsize
\conceptlead{The code mirrors the equation. Shapes: batch $B$, length $T$, key width $d_k$.}
{\scriptsize\begin{verbatim}
import torch
import torch.nn.functional as F

# X: [B,T,d_model], Wq/Wk/Wv: [d_model,d_k]
T = X.size(1)
Q, K, V = X @ Wq, X @ Wk, X @ Wv
scores = (Q @ K.transpose(-2, -1)) / Q.size(-1)**0.5
mask = torch.triu(torch.ones(T,T,device=X.device,
                   dtype=torch.bool), diagonal=1)
scores = scores.masked_fill(mask, float('-inf'))
weights = F.softmax(scores, dim=-1)
out = weights @ V
\end{verbatim}}
\begin{keyidea}A full model adds padding masks, multi-head projections, normalization, MLPs, and training code.\end{keyidea}
\end{frame}

\section{What Transformers Fix---and What They Do Not}

\begin{frame}{Full Attention Has a Quadratic Score Matrix}
\small
\begin{columns}[T,onlytextwidth]
\begin{column}{.59\textwidth}\lecturevisual[2.57in]{quadratic_attention.png}\end{column}
\begin{column}{.38\textwidth}
\begin{itemize}
\item Standard full self-attention forms $n\times n$ pairwise scores for a length-$n$ sequence.
\item At $n=4096$, that is about 16.8 million scores \emph{per head}, before considering batching and layers.
\item Memory-efficient kernels and alternative attention designs change practical costs.
\end{itemize}
\end{column}
\end{columns}
\end{frame}

\begin{frame}{A Transformer Can Use Context Without Understanding Truth}
\small
\conceptlead{Attention improves the path by which information moves, not the reliability of every conclusion.}
\begin{itemize}
\item \textbf{Context window:} a model cannot directly attend to text excluded from its available context.
\item \textbf{Training mismatch:} predicting likely text is different from checking external facts.
\item \textbf{Data bias:} corpus patterns, including mistakes and stereotypes, may be learned.
\item \textbf{Interpretability:} an attention heatmap alone is not a complete causal explanation of an output.
\end{itemize}
\begin{keyidea}Architectural improvements do not remove the need to evaluate outputs on the task that matters.\end{keyidea}
\end{frame}

\begin{frame}{Putting the Whole Path Back Together}
\small
\lecturevisual[2.48in]{llm_pipeline.png}
\begin{enumerate}
\item \textbf{Represent} text as tokens, IDs, embeddings, and positional information.
\item \textbf{Contextualize} with Q/K/V attention, multiple heads, and stacked blocks.
\item \textbf{Predict} a distribution over the next token using a language-model head.
\item \textbf{Train} with masked next-token loss; \textbf{generate} by iterating one token at a time.
\end{enumerate}
\end{frame}

\begin{frame}{Exit Ticket: Explain the Architecture}
\small
\begin{enumerate}
\item For the sentence \emph{``The robot saw the charging station''}, where do Q, K, and V come from?
\item If a query receives attention weights $[0.2,0.3,0.5]$, what do we do with its three value vectors?
\item Why does a decoder-only generator need a causal mask when training?
\item What is the difference between encoder--decoder translation and decoder-only next-token generation?
\item What important limitation remains even if attention learns very long-range relationships?
\end{enumerate}
\begin{keyidea}If you can explain both the \emph{mathematics inside a head} and the \emph{dataflow across blocks}, you can reason about modern transformer language models.\end{keyidea}
\end{frame}

\begin{frame}{References and Source Notes}
\scriptsize
\begin{itemize}
\item Primary instructional sources: the two supplied lecture PDFs, \emph{Lecture 24: LLM pipeline} (10 pages) and \emph{Lecture 24: LLM -- Transformer} (29 pages).
\item Bahdanau, Cho, Bengio (2015), \emph{Neural Machine Translation by Jointly Learning to Align and Translate} (arXiv 2014).
\item Vaswani et al. (2017), \emph{Attention Is All You Need}, NeurIPS 2017. \url{https://arxiv.org/abs/1706.03762}.
\item Devlin et al. (2019), \emph{BERT: Pre-training of Deep Bidirectional Transformers for Language Understanding}.
\item Radford et al. (2018), \emph{Improving Language Understanding by Generative Pre-Training}; Brown et al. (2020), \emph{Language Models are Few-Shot Learners}.
\item Figures and exact toy calculations generated locally by \texttt{generate-figures.py}, with no trained model inference. Any language-dependency figure is \emph{schematic}, not a measured attention attribution.
\item Pedagogical correction to the source decks: the 2017 model was encoder--decoder; decoder-only next-token training is a related architecture; generated examples and sample token splits are illustrations, not specific tokenizer outputs.
\end{itemize}
\end{frame}
\end{document}
